% 第12章（本部分第5章）图：前代与回代的计算顺序与依赖，四阶三角矩阵。
% 左：前代计算 y_3 时用到 l_{3,1},l_{3,2} 与已算出的 y_1,y_2；右：回代计算 x_2 时用到 u_{2,3},u_{2,4} 与已算出的 x_3,x_4。
\begin{tikzpicture}[x=1mm,y=1mm,>=Stealth,
  cell/.style={rectangle,draw=FigureStroke,fill=white,line width=.5pt,minimum width=7mm,minimum height=7mm,inner sep=0pt,font=\figurefont\fontsize{7.5}{9}\selectfont},
  used/.style={cell,fill=FigureBlueFill,draw=FigureBlue},
  piv/.style={cell,fill=FigureYellowFill,draw=FigureYellow},
  done/.style={cell,fill=FigureGreenFill,draw=FigureGreen},
  note/.style={font=\FigureNoteFont\fontsize{7.5}{9}\selectfont,text=FigureMuted}]
\def\s{7}
% ---------- left panel: forward substitution, row i=3 ----------
\node[font=\figurefont\fontsize{8}{10}\selectfont] at (14,4.5) {前代 $\bm L\bm y=\bm c$};
\foreach \i in {1,...,4}{\foreach \j in {1,...,4}{
  \pgfmathtruncatemacro{\ok}{\j<=\i}
  \ifnum\ok=1
    \node[cell] (L\i\j) at ({(\j-.5)*\s},{-(\i-.5)*\s}) {};
  \fi
}}
\node[used] at (L31) {$l_{3,1}$};
\node[used] at (L32) {$l_{3,2}$};
\node[piv]  at (L33) {$l_{3,3}$};
\foreach \i/\st/\lab in {1/done/{y_1},2/done/{y_2},3/piv/{y_3},4/cell/{y_4}}{
  \node[\st] (y\i) at (35.5,{-(\i-.5)*\s}) {$\lab$};
}
\draw[->,FigureStroke,line width=.8pt] (42,-1.5)--(42,-26.5);
\node[note,anchor=west,align=left] at (43,-14) {$i=1$\\至$4$};
% ---------- right panel: back substitution, row i=2 ----------
\begin{scope}[xshift=57mm]
\node[font=\figurefont\fontsize{8}{10}\selectfont] at (14,4.5) {回代 $\bm U\bm x=\bm y$};
\foreach \i in {1,...,4}{\foreach \j in {1,...,4}{
  \pgfmathtruncatemacro{\ok}{\j>=\i}
  \ifnum\ok=1
    \node[cell] (U\i\j) at ({(\j-.5)*\s},{-(\i-.5)*\s}) {};
  \fi
}}
\node[piv]  at (U22) {$u_{2,2}$};
\node[used] at (U23) {$u_{2,3}$};
\node[used] at (U24) {$u_{2,4}$};
\foreach \i/\st/\lab in {1/cell/{x_1},2/piv/{x_2},3/done/{x_3},4/done/{x_4}}{
  \node[\st] (x\i) at (35.5,{-(\i-.5)*\s}) {$\lab$};
}
\draw[->,FigureStroke,line width=.8pt] (42,-26.5)--(42,-1.5);
\node[note,anchor=west,align=left] at (43,-14) {$i=4$\\至$1$};
\end{scope}
% ---------- notes ----------
\node[note] at (22,-33) {$y_3=(c_3-l_{3,1}y_1-l_{3,2}y_2)/l_{3,3}$};
\node[note] at (79,-33) {$x_2=(y_2-u_{2,3}x_3-u_{2,4}x_4)/u_{2,2}$};
\end{tikzpicture}
